%%==========================================================================%%
%% abstract.tex -- working abstract for the pediatric whole-body FDG-PET/MRI  %%
%% lymphoma lesion dataset descriptor (Nature Scientific Data).              %%
%%                                                                           %%
%% Constraints: 170 words maximum, no references, no URLs or data-access     %%
%% details, no sub-headings, no claims of new scientific findings.           %%
%%                                                                           %%
%% NOTE: Springer Nature requires single-file submission. Before submitting, %%
%% inline the contents of this file into main.tex.                           %%
%%==========================================================================%%

\abstract{Whole-body FDG-PET combined with MRI is established for staging and treatment monitoring in pediatric lymphoma, where it limits the cumulative radiation of repeated examinations.
Public datasets that drive automated PET lesion analysis are almost entirely adult and computed-tomography-based, and models built on them transfer poorly to children.
We present a single-center retrospective cohort of pediatric and young-adult lymphoma patients imaged with whole-body [$^{18}$F]fluorodeoxyglucose PET and paired MRI.
The release provides the PET and paired MRI series, manually drawn tumor-lesion segmentations on the FDG-avid examinations, disease-free follow-up examinations as pathology-negative references, and indexing metadata.
We describe the cohort, the PET and MRI acquisition, the de-identification and annotation procedures, the file organization, and a technical validation of label quality.
The dataset is intended for the development and benchmarking of automated PET-guided lesion detection and segmentation, tumor-burden quantification, and adult-to-pediatric transfer learning.}
